% Exercises the manual citation-node commands, \usescite (inside a proof) and
% \proofgraphciteedge (anywhere), including a result that has no proof at all
% to hold a citation. The cite option is deliberately *not* given: an explicit
% command draws its citation node regardless, while a plain \cite still records
% nothing. A key is cited manually twice with different notes, checking that the
% (key, note) slots are kept apart, and t2's list form checks that a
% comma-separated argument yields one node per key.
% citelabel=key keeps the fixture bibliography-style independent, and the
% manual \bibitem keeps it bibtex-free.
\documentclass{article}
\usepackage{amsthm}
\usepackage[citelabel=key]{proofgraph}

\newtheorem{theorem}{Theorem}

\begin{document}

\begin{theorem}\label{t1}One.\end{theorem}
\begin{proof}
  Known, see \cite{bk}, which the cite option would have captured.%
  \usescite[Thm~3]{bk}%
\end{proof}

\begin{theorem}\label{t2}Two, stated without a proof.\end{theorem}

\proofgraphciteedge{t2}{bk, art}
\proofgraphciteedge[Lem~1]{t1}{art}

\begin{thebibliography}{9}
\bibitem{bk}A book.
\bibitem{art}An article.
\end{thebibliography}
\end{document}
